\documentclass[aspectratio=169, handout]{beamer}

%%%%%%%%%%%%%%%%%%%%%%%%%%%%%%%%%%%%%%%%%%%%%%%%%%%%%%%%%%%%%%%%%%%%%%%%%%%%%%%%%%%%%%%%%%%%%%%%%%%%%%%%%%%%%%%%%%%%%%%
%
%                                                      PREAMBLE.STY
%
% Elements configured here:
%
%       -> Control panel (high level control for LANGUAGE and COLORTHEME)
%       -> Language Support
%       -> Colors (switches between different color themes) [/CodexStyle/colorthemes/...]
%       -> General Packages
%       -> Document Layout (switches between different layouts) [/CodexStyle/layouts/...]
%       -> Math Setup (math packages)
%       -> Fonts (Body)
%       -> Fonts (Math)
%       -> Tables
%       -> Graphical Setup (tikz pgfplots)
%
%%%%%%%%%%%%%%%%%%%%%%%%%%%%%%%%%%%%%%%%%%%%%%%%%%%%%%%%%%%%%%%%%%%%%%%%%%%%%%%%%%%%%%%%%%%%%%%%%%%%%%%%%%%%%%%%%%%%%%%

%!%!%!%!%!%!%!%!%!%!%!%!%!%!%!%!%!%!%!%!%!%!%!%!%
%                 Control Panel                 %
%!%!%!%!%!%!%!%!%!%!%!%!%!%!%!%!%!%!%!%!%!%!%!%!%


\newcommand{\CodexLanguage}{french} % set language
\newcommand{\CodexColorTheme}{CodexStyle/colorthemes/colors_marine} % set colortheme (path to desired .sty)


%%%%%%%%%%%%%%%%%%%%%%%%%%%%%%%%%%%%%%%%%%%%%%%%%
%                Language Support               %
%%%%%%%%%%%%%%%%%%%%%%%%%%%%%%%%%%%%%%%%%%%%%%%%%

\usepackage[\CodexLanguage]{babel} % language support priority language last (set in control panel)

% French renamings ------------------------------
\addto\captionsfrench{\renewcommand\frenchtablename{{\FBfigtabshape Tableau}}} % rename table name (french)
% \addto\captionsfrench{%
%   \renewcommand{\refname}{Bibliographie}%
% }


%%%%%%%%%%%%%%%%%%%%%%%%%%%%%%%%%%%%%%%%%%%%%%%%%
%                     Colors                    %
%%%%%%%%%%%%%%%%%%%%%%%%%%%%%%%%%%%%%%%%%%%%%%%%%

\usepackage[x11names]{xcolor} % adds in typical color features
\usepackage{\CodexColorTheme} % choose a style (set in control panel)


%%%%%%%%%%%%%%%%%%%%%%%%%%%%%%%%%%%%%%%%%%%%%%%%%
%               General Packages                %
%%%%%%%%%%%%%%%%%%%%%%%%%%%%%%%%%%%%%%%%%%%%%%%%%

\usepackage{setspace}  % line spacing
% \usepackage{microtype} % small typography tweaks
\usepackage{graphicx}  % figure insertion
\usepackage{subfig}	   % allows for subfigures MIGRATE TO A MODERN PACKAGE MAYBE
\usepackage{geometry}  % document geometry
\usepackage{fancybox}  % boxes, frames etc
\usepackage{url}       % typographically sound URLs
\usepackage{float}     % Fixes hacky geometry problems
\usepackage[
format    = hang,
margin    = 5mm,
font      = small,
labelfont = bf,
labelsep  = space
]{caption}             % tweak caption layout and format
\usepackage{csquotes}  % For quotes
\usepackage[
    backend  = biber,
    style    = nature,
    sorting  = none,
    backref  = true
]{biblatex}
\usepackage{hyperref}[ % hyperref with some setup
    hyperindex  = true,
    bookmarks   = true,
    pdfa
]
% more hyperref options
\hypersetup{
    colorlinks=true,         % colorize links
    breaklinks=true,         % carriage return for long links
    urlcolor=URLColor,       % hyperlink url color
    linkcolor=LinkColor,     % internal link color
    citecolor=CiteColor,     % citation link color
    bookmarksopen=true,      % opens pdf bookmarks in the beginning
}


%%%%%%%%%%%%%%%%%%%%%%%%%%%%%%%%%%%%%%%%%%%%%%%%%
%                 Document Layout               %
%%%%%%%%%%%%%%%%%%%%%%%%%%%%%%%%%%%%%%%%%%%%%%%%%

% Switches between longform and shortform automatically
\makeatletter
\@ifclassloaded{book}{ %%% BOOK %%%
    \usepackage{CodexStyle/layouts/longform_layout} % load longform
}{
\@ifclassloaded{article}{ %%% ARTICLE %%%
    \usepackage{CodexStyle/layouts/shortform_layout} % load shortform
}{
\@ifclassloaded{beamer}{ %%% BEAMER %%%
    \usepackage{CodexStyle/layouts/presentation_layout} % load presentation
}{
    \PackageError{mydoc}
        {Unsupported document class!}
        {`book` -> longform, `article` -> shortform, `beamer` -> presentation}
}}}
\makeatother


%%%%%%%%%%%%%%%%%%%%%%%%%%%%%%%%%%%%%%%%%%%%%%%%%
%                   Math Setup                  %
%%%%%%%%%%%%%%%%%%%%%%%%%%%%%%%%%%%%%%%%%%%%%%%%%

\usepackage{amsmath}          % math environments
\usepackage{mathtools}        % tools for math formating
\usepackage[nice]{nicefrac}   % nicer fractions
\usepackage{cancel}           % allows to scratch expressions.
\usepackage{slashed}          % allows to slash individual characters.
\usepackage{xargs}            % better handling of optional arguments for commands
\usepackage{empheq}           % colored boxes in math env
\usepackage[most]{tcolorbox}  % adds colorful box environments

%%% Import custom math macros
\usepackage{CodexStyle/macros} % custom latex macros (macros.sty)


%%%%%%%%%%%%%%%%%%%%%%%%%%%%%%%%%%%%%%%%%%%%%%%%%
%                  Fonts (Body)                 %
%%%%%%%%%%%%%%%%%%%%%%%%%%%%%%%%%%%%%%%%%%%%%%%%%

\usepackage{fontspec} % Finer font selection (requires Lua/XeLaTeX)
\usepackage{unicode-math}[ % Finer math font selection
]

% Main font switches based on documentclass
\makeatletter
\@ifclassloaded{book}{ %%% BOOK %%%
    \setmainfont{EB Garamond}[ % EB Garamond
        Ligatures   = TeX, % ligatures
        OpticalSize = On, % optical weight adjustment
        % Numbers     = OldStyle, % can be a little hard to read
        SmallCapsFeatures = {LetterSpace=5} % Adds spaces for small caps
    ]
}{
\@ifclassloaded{article}{ %%% ARTICLE %%%
    % \setmainfont{Libertinus Serif}[ % Libertinus in short documents
    %     Ligatures = TeX,
    %     Kerning   = On,
    % ]
    \setmainfont{XCharter}[ % XCharter in short documents
    ]
}{
\@ifclassloaded{beamer}{ %%% BEAMER %%%
    \setmainfont{XCharter}[ % XCharter in short documents
    ]
    \setsansfont{Open Sans}
}{ %%% ELSE %%%
    \PackageError{mydoc} % throws an error
        {Unsupported document class!}
        {`book` -> longform, `article` -> shortform, `beamer` -> presentation}
}}}
\makeatother


%%%%%%%%%%%%%%%%%%%%%%%%%%%%%%%%%%%%%%%%%%%%%%%%%
%                  Fonts (Math)                 %
%%%%%%%%%%%%%%%%%%%%%%%%%%%%%%%%%%%%%%%%%%%%%%%%%

\setmathfont{XCharter Math}[ % nice Charter (near) universal math font
    % Scale=MatchLowercase % set scale
]
\setmathfont{Latin Modern Math}[ % A few swaps and additions
    range = {\dagger, \ddagger, \mathcal}, % less dramatic daggers and distinct mathcal
    % Scale=MatchLowercase % set scale
]


%%%%%%%%%%%%%%%%%%%%%%%%%%%%%%%%%%%%%%%%%%%%%%%%%
%                     Tables                    %
%%%%%%%%%%%%%%%%%%%%%%%%%%%%%%%%%%%%%%%%%%%%%%%%%

\usepackage{array} % tabular functions
\usepackage{dcolumn} % allows for aligning of values wrt to the decimal
\newcolumntype{C}{>{$\displaystyle}c<{$}} % centered math column
\newcolumntype{L}{>{$\displaystyle}l<{$}} % left aligned math column
\newcolumntype{R}{>{$\displaystyle}r<{$}} % right aligned math column
\renewcommand{\arraystretch}{1.5} % vertical spacing of tables


%%% UNDER CONSTRUCTION %%%
%-----------------------------------------------
%%%%%%%%%%%%%%%%%%%%%%%%%%%%%%%%%%%%%%%%%%%%%%%%%
%                 Graphical Setup               %
%%%%%%%%%%%%%%%%%%%%%%%%%%%%%%%%%%%%%%%%%%%%%%%%%

\usepackage{tikz}
\usetikzlibrary{
	calc,
	patterns,
	positioning,
    external,
    shapes,
    fit,
    backgrounds,
    arrows.meta,
	positioning,
    decorations,
    decorations.pathmorphing,
    decorations.markings,
    shapes.geometric,
    arrows
}
% \tikzexternalize[prefix=figures/pdf/]

\usepackage{pgfplots}
\pgfplotsset{compat=1.16}
\pgfdeclarelayer{background}
\pgfdeclarelayer{foreground}
\pgfsetlayers{background,main,foreground}

\usepackage{CodexStyle/layouts/tikzstyles}
%%%%%%%%%%%%%%%%%%%%%%%%%%%%%%%%%%%%%%%%%%%%%%%%


\title[Geometry in QM]{Geometry in Quantum Theory}
\subtitle{...with a dash of condensed matter}
\author{Presented by\\Théo N. Dionne}

\institute[]{
    Institut Quantique \\
    Sherbrooke, Québec \\
    October 22 2026
}

\date{} % suppress date field

\titlegraphic{
    \includegraphics[width=3cm]{./figures/logos/UdeS_logo_h_noir_cropped.pdf}
    \hspace{1cm}
    \includegraphics[width=4cm]{./figures/logos/IQ_logo_h_noir_cropped.pdf}
}

\addbibresource{sample.bib}


\usepackage[percent]{overpic}
\begin{document}

\begingroup\setbeamertemplate{navigation symbols}{}

\begin{frame}[noframenumbering]
	\titlepage
\end{frame}


% \begin{frame}[noframenumbering]\frametitle{Table of contents}
%     \tableofcontents
% \end{frame}

\endgroup

\begin{frame}\frametitle{Geometry and Physics}
\begin{columns}
	\column{0.6\linewidth}
	    \visible<2->{\begin{columns} % Geometric optics
			\column{0.5\linewidth}\centering
			    {\Large Geometrical Optics}
			\column{0.5\linewidth}\centering
			\begin{overpic}[width=\linewidth]{figures/pdf/Lens3b.pdf}
                \put (55,3) {\tiny [Wikipedia, Optics]}
            \end{overpic}
		\end{columns}}
        \visible<3->{\begin{columns} % General Relativity
            \column{0.5\linewidth}\centering
                {\Large General Relativity}
            \column{0.5\linewidth}\centering
            \begin{overpic}[width=\linewidth]{figures/pdf/Spacetime_lattice_analogy_white-1.pdf}
                \put (33,3) {\tiny [Wikipedia, General Relativity]}
            \end{overpic}
        \end{columns}}
        \visible<4->{\begin{columns} % Quantum Field Theory
            \column{0.5\linewidth}\centering
                {\Large Quantum Field Theory}
            \column{0.5\linewidth}\centering
                \begin{overpic}[width=\linewidth]{figures/GaugeFixing.png}
                    \put (45,63) {\tiny [Fradkin, QFT Integrated]}
                \end{overpic}
        \end{columns}}
	\column{0.4\linewidth}\centering
	\visible<5->{How about Quantum Mechanics?}

	\visible<6->{Is Quantum Geometry$^{\text{TM}}$ \textbf{new?}
	\begin{overpic}[width=0.6\linewidth]{./figures/QuantumGeometryTrend.png}
	    \put(5,55){\tiny[Google Trends]}
	\end{overpic}}

	\visible<7->{\textbf{No!} Geometry is \textbf{\color{MainColor}foundational}\vspace{1ex}

	\begin{overpic}[width=0.7\linewidth]{figures/foundations.jpg}
        \put(34,30){\bfseries\Large\color{MainColor}1932}
        \put(18,2){\tiny\color{white}[Maison Tajan]}
    \end{overpic}}

\end{columns}


\end{frame}

\begin{frame}\frametitle{Roadmap}
    \center{\LARGE\sffamily\bfseries\color{HighlightColor}
        \enquote{Answer me these questions three...}
    }
    \vspace{2ex}
    {\Large\begin{enumerate}
        \item<2-> \textbf{\color{MainColor}What} is the \textbf{shape} of \textbf{inequivalent} quantum \textbf{states}?
        \item<3-> \textbf{\color{MainColor}How} do we \textbf{quantify geometry}?
        \item<4-> \textbf{\color{MainColor}Where} is \textbf{geometry} in \textbf{materials}?
    \end{enumerate}}
\end{frame}

\begingroup\setbeamertemplate{navigation symbols}{}
\begin{frame}[noframenumbering]
    \center{
	    \LARGE\sffamily\bfseries\color{HighlightColor} What do quantum states look like?
	}
	{\center{\resizebox{0.3\linewidth}{!}{\begin{tikzpicture}

\node[above] at (2,3) {\footnotesize$S^{2n-1}\subset\mathbb{C}^{n}$}; % annotate bundle

% Draw fibre bundle bottom
\draw[thick] (0,0) -- (2,0);
\draw[thick] (2,0) -- (3,1);
\draw[thick] (3,1) -- (1,1);
\draw[thick] (1,1) -- (0,0);

% Draw projective space section
\fill[MainColor] (0,0) -- (2,0) -- (3,1) -- (1,1) -- cycle;
\node[right] at (2,0) {\footnotesize\color{MainColor}$\mathbb{CP}^{n-1}$}; % annotate section

% Draw ray
\draw[thick, HighlightColor] (1.25,0.75) -- (1.25,2+0.75);
\filldraw[fill=white, draw=HighlightColor] (1.25,0.75) circle (0.04);
\fill[HighlightColor] (1.25,2+0.75) circle (0.04);
\node[below] at (1.25,0.75) {\footnotesize\color{white}$[\psi]$}; % annotate ray

% Draw point on ray
\fill (1.25,1.35) circle (0.04);
\node[right] at (1.25,1.35) {\footnotesize$\ket{\psi}$};
\fill (1.25,2.5) circle (0.04);
\node[right] at (1.25,2.5) {\footnotesize$\Exp{i\phi}\!\ket{\psi}$};

% Draw fibre bundle top
\draw[thick] (0,0+2) -- (2,0+2);
\draw[thick] (2,0+2) -- (3,1+2);
\draw[thick] (3,1+2) -- (1,1+2);
\draw[thick] (1,1+2) -- (0,0+2);

% Draw fibre bundle sides
\draw[thick] (0,0) -- (0,0+2);
\draw[thick] (2,0) -- (2,0+2);
\draw[thick] (3,1) -- (3,1+2);
\draw[thick] (1,1) -- (1,1+2);

% \draw[MainColor, -stealth, thick] (0,0) -- (2,0);
% \draw[MainColor, -stealth, thick] (2,0) -- node[midway, right] {\color{MainColor}$[\psi]$} (3,1);
% \draw[MainColor, -stealth, thick] (3,1) -- (1,1);
% \draw[MainColor, -stealth, thick] (1,1) -- (0,0);

% \draw[dashed] (0,0) -- (0,2);
% \draw[dashed] (2,0) -- (2,2);
% \draw[dashed] (3,1) -- (3,3);
% \draw[dashed] (1,1) -- (1,3);

% \draw[black, -stealth, thick] (0,0) -- (2,0+0.3);
% \draw[black, -stealth, thick] (2,0+0.3) -- (3,1+2*0.3);
% \draw[black, -stealth, thick] (3,1+2*0.3) -- node[near start, above] {$\psi$} (1,1+3*0.3);
% \draw[black, -stealth, thick] (1,1+3*0.3) -- (0,0+4*0.3);

% \draw[HighlightColor, thick] (0,0) -- node[midway, left] {\color{HighlightColor}$\Exp{i\phi}$} (0,1.2);

% \fill[black] (0,0) circle (0.03);
% \fill[black] (0,1.2) circle (0.03);

\end{tikzpicture}
}}}
\end{frame}
\endgroup

\begin{frame}\frametitle{Quantum Theory}
\begin{columns}
	\column{0.8\linewidth}
	\visible<2->{\textbf{Theories such that...}}
    \begin{itemize}
        \item<3->[\bullet] \textbf{\color{MainColor}State:} Complex vector $\ket{\psi}\in\mathbb{C}^{n}$
        \item<4->[\bullet] \textbf{\color{MainColor}Observable:} Self-adjoint operators $\op{A}:\mathbb{C}^{n}\to\mathbb{C}^{n}~|~\op{A}^\dagger = \op{A}$
        \item<5->[\bullet] \textbf{\color{MainColor}Expectation Value:} \color{black} Real number \alert<7->{$\avg{\op{A}}\equiv\frac{\bra{\psi}\op{A}\ket{\psi}}{\braket{\psi}{\psi}}\in\mathbb{R}$\uncover<7->{(!)}}
    \end{itemize}
    \vspace{2ex}
    \uncover<6->{
    \inlinecbox{
        {Relevant geometry of quantum states is $\symit{\mathbb{C}^n}$?... }\uncover<7->{\alert{Not quite!}}
    }
    }
	\column{0.2\linewidth}
	\visible<8->{\textbf{\color{MainColor}Norm}\\$\ket{\psi}\sim{\color{MainColor}\alpha}\ket{\psi}$\\\vspace{10ex}}
	\visible<9->{\textbf{\color{HighlightColor}Phase}\\$\ket{\psi}\sim\alert{\Exp{i\phi}}\ket{\psi}$}
\end{columns}

\end{frame}

\begin{frame}\frametitle{Norm Redundancy}
\begin{columns}
	\column{0.6\linewidth}
        \visible<2->{For convenience, normalize to unity $\left|\ket{\psi}\right|=1$
        \eqn{
            \ket{\psi} \to \frac{\ket{\psi}}{\sqrt{\braket{\psi}{\psi}}}
        }}
        \visible<3->{Normalized states $\ket{\psi}\in\mathbb{C}^{n}$ define hypersphere
        \eqn{
            \braket{\psi}{\psi} = \sum_{i=1}^{n}\abs{c_i}^2 = \underbrace{\sum_{i=1}^{n}[(c_i')^2+(c_i'')^2]=1}_{2n\text{ real variables}}
        }}
        \visible<4->{\textbf{Distinct} states in $\symbf{S^{2n-1}}\subset\mathbb{C}^n$}
	\column{0.4\linewidth}\centering
	    \visible<3->{\resizebox{\linewidth}{!}{\begin{tikzpicture}

% Axes just cause
\draw[thick, -Stealth] (0,0) -- node[at end, right, xshift=-0.cm] {\large$\mathbb{C}^{n}$} (0,4);
\draw[thick, -Stealth] (0,0) -- (4,0);

% Circle to guide the eye
\draw[MainColor, thick] (2,2) circle (1.5cm);
\node at (3.4, 3.6) {\color{MainColor}$\left|\ket{\psi}\right|=1$};
\node at (2,2) {\large{\color{MainColor}$S^{2n-1}$}};

\end{tikzpicture}
}}
	    \visible<5->{\textbf{\alert{What about phase?}}}
\end{columns}

\end{frame}

\begin{frame}\frametitle{Phase Redundancy}
\begin{columns}
	\column{0.6\linewidth}
	\visible<2->{Pull out one factor of phase
	\eqn{
	    \ket{\psi} \sim \Exp{i\phi}\!\ket{\psi}
	}}
	\visible<3->{For example, factor out $\Exp{i\phi_1}$
    \eqn{
        &\pmat{r_1\Exp{i\phi_1}, r_2\Exp{i\phi_2},\dots,r_n\Exp{i\phi_n}}^{\symup{T}}\\
        \sim\Exp{i\phi_1}&\pmat{{\color{MainColor}r_1}, {\color{MainColor}r_2}{\color{HighlightColor}\Exp{i(\phi_2-\phi_1)}},\dots,{\color{MainColor}r_n}{\color{HighlightColor}\Exp{i(\phi_n-\phi_1)}}}^{\symup{T}}
    }}
    \visible<4->{Under normalization, all distinct states
    \eqn{
        2n-1 \overset{\text{Normalization}}{\longrightarrow} 2n-2~\text{(real dimension)}
    }}

	\column{0.4\linewidth}
	\visible<3->{\center{
	    \resizebox{\linewidth}{!}{\begin{tikzpicture}

% Axes just cause
\draw[thick, -Stealth] (0,0) -- (0,4);
\node[below] at (4,0) {$r_1$};
\draw[thick, -Stealth] (0,0) -- (4,0);
\node[left] at (0,4) {$r_2$};

% Add a tick
\node[below] at (3,0) {$1$};
\node[left] at (0,3) {$1$};

% Draw circular shell
\draw[MainColor, thick] (3,0) arc[start angle=0, delta angle=90, radius=3cm] node[pos=0.5, circle, minimum size=0.1cm, inner sep=0, fill=black, name=point] {};
\node at (1.5,1.5) {\Large\color{MainColor}Norms}; % annotate shell
\node at (1.5,1.1) {\Large\color{MainColor}}; % annotate shell

% Draw U(1) circle
\draw[HighlightColor, thick] (3,3) circle (0.5cm);
\node at (3,3) {\color{HighlightColor}$U\!(1)$};
\node at (4.1,3.5) {\color{HighlightColor}$\times(n-1)$};

% Draw arrow
\draw[-stealth, thick] (2.8,2.5) to[bend left] (point);

\end{tikzpicture}
}
	}}
	\visible<5->{\textbf{What does this look like?}}

\end{columns}
\end{frame}

\begin{frame}\frametitle{Example: Bloch Sphere}
\begin{columns}
	\column{0.5\linewidth}
	    \visible<2->{\textbf{Two level} system $\ket{\psi}\in\mathbb{C}^{2}$
					\eqn{\operatorname{dim}_{\mathbb{R}}(\psi)=4~~\overset{2n-2}{\Longrightarrow}~~ \symbf{2}~\text{distinct \textbf{dimensions}}}}
		\visible<3->{For $u,v\in\mathbb{R}^{+}$ and $\Exp{i\phi}\in U(1)$
		\eqn{\ket{\psi}=\pmat{{\color{MainColor}u},{\color{MainColor}v}{\color{HighlightColor}\Exp{i\phi}}}^{\symup{T}}}}
		\visible<4->{Trouble at $u=0$, no biggie}
	\column{0.5\linewidth}
        \visible<3->{\center{\resizebox{\linewidth}{!}{\begin{tikzpicture}

% draw phases
\draw[HighlightColor, thick] (0,0) ellipse (2 and 0.15);
\draw[HighlightColor, thick] (0,1.414) ellipse (1.414 and 0.15);
\draw[HighlightColor, thick] (0,-1.414) ellipse (1.414 and 0.15);

% draw half circles, one full one dashed
\draw[MainColor,thick] (0,2) arc[radius=2cm, start angle=90, delta angle=-180];
\draw[MainColor,thick, dashed] (0,2) arc[radius=2cm, start angle=90, delta angle=180];

% Nodes for remarkable points
\node[name=north, circle, inner sep=0, minimum size=0.1cm, fill=black] at (0,2) {};
\node[above] at (north) {$(1,0)^{\symup{T}}$};
\node[name=east, circle, inner sep=0, minimum size=0.1cm, fill=black] at (2,0) {};
\node[right] at (east) {$\frac{1}{\sqrt{2}}(1,1)^{\symup{T}}$};
\node[name=south, circle, inner sep=0, minimum size=0.1cm, fill=black] at (0,-2) {};
\node[below] at (south) {$(0,1)^{\symup{T}}$};
\node[name=green, circle, inner sep=0, minimum size=0.1cm, fill=black] at (0,-0.15) {};
\node[below] at (green) {$\frac{1}{\sqrt{2}}(1,i)^{\symup{T}}$};
\node[name=west, circle, inner sep=0, minimum size=0.1cm, fill=black] at (-2,0) {};
\node[left] at (west) {$\frac{1}{\sqrt{2}}(1,-1)^{\symup{T}}$};

\end{tikzpicture}
}}}
		\visible<5->{More \textbf{complicated} for $\symbf{n>2}$}
\end{columns}
\end{frame}

\begin{frame}\frametitle{Complex Projective Space}
\begin{columns}
	\column{0.6\linewidth}
	\visible<2->{\textbf{Inequivalent} states: \textbf{complex projective} space
	\eqn{
	    \mathbb{C}^{n}&{\color{MainColor}\overset{\text{normalize}}{\longrightarrow}}S^{2n-1}{\color{HighlightColor}\overset{\symup{mod}~U(1)}{\longrightarrow}}\mathbb{CP}^{n-1}
	}}
	\visible<3->{$[\psi]$ in $\mathbb{CP}^{n-1}$ = \textbf{ray} of equivalent states
	\eqn{
	    \mathbb{C}^{n} ~~\overset{\ket{\psi}\sim z\ket{\psi}}{\longrightarrow}~~\mathbb{CP}^{n-1}
	}}
	\column{0.4\linewidth}
	\visible<3->{\resizebox{0.8\linewidth}{!}{\begin{tikzpicture}

\node[above] at (2,3) {\footnotesize$S^{2n-1}\subset\mathbb{C}^{n}$}; % annotate bundle

% Draw fibre bundle bottom
\draw[thick] (0,0) -- (2,0);
\draw[thick] (2,0) -- (3,1);
\draw[thick] (3,1) -- (1,1);
\draw[thick] (1,1) -- (0,0);

% Draw projective space section
\fill[MainColor] (0,0) -- (2,0) -- (3,1) -- (1,1) -- cycle;
\node[right] at (2,0) {\footnotesize\color{MainColor}$\mathbb{CP}^{n-1}$}; % annotate section

% Draw ray
\draw[thick, HighlightColor] (1.25,0.75) -- (1.25,2+0.75);
\filldraw[fill=white, draw=HighlightColor] (1.25,0.75) circle (0.04);
\fill[HighlightColor] (1.25,2+0.75) circle (0.04);
\node[below] at (1.25,0.75) {\footnotesize\color{white}$[\psi]$}; % annotate ray

% Draw point on ray
\fill (1.25,1.35) circle (0.04);
\node[right] at (1.25,1.35) {\footnotesize$\ket{\psi}$};
\fill (1.25,2.5) circle (0.04);
\node[right] at (1.25,2.5) {\footnotesize$\Exp{i\phi}\!\ket{\psi}$};

% Draw fibre bundle top
\draw[thick] (0,0+2) -- (2,0+2);
\draw[thick] (2,0+2) -- (3,1+2);
\draw[thick] (3,1+2) -- (1,1+2);
\draw[thick] (1,1+2) -- (0,0+2);

% Draw fibre bundle sides
\draw[thick] (0,0) -- (0,0+2);
\draw[thick] (2,0) -- (2,0+2);
\draw[thick] (3,1) -- (3,1+2);
\draw[thick] (1,1) -- (1,1+2);

% \draw[MainColor, -stealth, thick] (0,0) -- (2,0);
% \draw[MainColor, -stealth, thick] (2,0) -- node[midway, right] {\color{MainColor}$[\psi]$} (3,1);
% \draw[MainColor, -stealth, thick] (3,1) -- (1,1);
% \draw[MainColor, -stealth, thick] (1,1) -- (0,0);

% \draw[dashed] (0,0) -- (0,2);
% \draw[dashed] (2,0) -- (2,2);
% \draw[dashed] (3,1) -- (3,3);
% \draw[dashed] (1,1) -- (1,3);

% \draw[black, -stealth, thick] (0,0) -- (2,0+0.3);
% \draw[black, -stealth, thick] (2,0+0.3) -- (3,1+2*0.3);
% \draw[black, -stealth, thick] (3,1+2*0.3) -- node[near start, above] {$\psi$} (1,1+3*0.3);
% \draw[black, -stealth, thick] (1,1+3*0.3) -- (0,0+4*0.3);

% \draw[HighlightColor, thick] (0,0) -- node[midway, left] {\color{HighlightColor}$\Exp{i\phi}$} (0,1.2);

% \fill[black] (0,0) circle (0.03);
% \fill[black] (0,1.2) circle (0.03);

\end{tikzpicture}
}}
	\visible<4->{\enquote{\textbf{Fibre bundle} approach}}
\end{columns}
\visible<5->{\center{\inlinecbox{\Large General yet impractical... how to keep track?}}}
\end{frame}


\begingroup\setbeamertemplate{navigation symbols}{}
\begin{frame}[noframenumbering]
	\center{
	    \LARGE\sffamily\bfseries\color{HighlightColor} How do we quantify state geometry?
	}

	{\center{
	    \resizebox{0.4\linewidth}{!}{\begin{tikzpicture}


% Draw projective space section
\fill[MainColor] (0,0) -- (2,0) -- (3,1) -- (1,1) -- cycle;
\node[right] at (2,0) {\footnotesize\color{MainColor}$\mathbb{CP}^{n-1}$}; % annotate section

% Draw ray
\draw[thick, HighlightColor] (1.25,0.75) -- (1.25,2+0.75);
\filldraw[fill=white, draw=HighlightColor] (1.25,0.75) circle (0.04);
\fill[HighlightColor] (1.25,2+0.75) circle (0.04);

% Draw neighboring ray
\draw[thick, HighlightColor] (1.25+0.7,0.75-0.2) -- (1.25+0.7,2+0.75-0.2);
\filldraw[fill=white, draw=HighlightColor] (1.25+0.7,0.75-0.2) circle (0.04);
\fill[HighlightColor] (1.25+0.7,2+0.75-0.2) circle (0.04);

% Draw path
\draw[thick, white] (1.25,0.75) to[bend right] node [below] {\color{white}$\D[\psi]^2$} (1.25+0.7,0.75-0.2);
\foreach \i in {1,2,3,4} { % draw suggestion of rays being compared
    \draw[thick, dashed, HighlightColor] (1.25,0.75+0.5*\i) to[bend right] (1.25+0.7,0.75-0.2+0.5*\i);
}

\end{tikzpicture}
}\resizebox{0.4\linewidth}{!}{\begin{tikzpicture}

\draw[MainColor, -stealth, thick] (0,0) -- (2,0);
\draw[MainColor, -stealth, thick] (2,0) -- node[midway, right] {\color{MainColor}$[\psi]$} (3,1);
\draw[MainColor, -stealth, thick] (3,1) -- (1,1);
\draw[MainColor, -stealth, thick] (1,1) -- (0,0);

\draw[dashed] (0,0) -- (0,2);
\draw[dashed] (2,0) -- (2,2);
\draw[dashed] (3,1) -- (3,3);
\draw[dashed] (1,1) -- (1,3);

\draw[black, -stealth, thick] (0,0) -- (2,0+0.3);
\draw[black, -stealth, thick] (2,0+0.3) -- (3,1+2*0.3);
\draw[black, -stealth, thick] (3,1+2*0.3) -- node[near start, above] {$\psi$} (1,1+3*0.3);
\draw[black, -stealth, thick] (1,1+3*0.3) -- (0,0+4*0.3);

\draw[HighlightColor, thick] (0,0) -- node[midway, left] {\color{HighlightColor}$\Exp{i\phi}$} (0,1.2);

\fill[black] (0,0) circle (0.03);
\fill[black] (0,1.2) circle (0.03);

\end{tikzpicture}
}
	}}
\end{frame}
\endgroup


\begin{frame}\frametitle{Metrics in a Nutshell}
    \begin{columns}
        \column{0.6\linewidth}
        \visible<2->{Parametrized vector curve $\ve{V}(\lambda)$ on manifold $\mathscr{M}$
        \eqn{
            \D\ve{V}^2 = \D\ve{V}\cdot\D\ve{V} &= \del_{\mu}\ve{V}\cdot\del_{\nu}\ve{V}\dd[x^{\mu}]{\lambda}\dd[x^{\nu}]{\lambda}\D\lambda^2\\
            &= {\color{MainColor}\del_{\mu}\ve{V}\cdot\del_{\nu}\ve{V}}\D x^{\mu}\D x^{\nu}
        }}
        \visible<3->{Arbitrary coordinates $\{x^{\mu}\}$
        \eqn{\color{MainColor}
            g_{\mu\nu} = \del_{\mu}\ve{V}\cdot\del_{\nu}\ve{V}
        }}
        \visible<4->{Metric {\color{MainColor}$g_{\mu\nu}$} encodes distances!}
        \column{0.4\linewidth}\centering
        \visible<2->{\resizebox{\linewidth}{!}{\begin{tikzpicture}

% \node[above] at (2,3) {\footnotesize$S^{2n-1}\subset\mathbb{C}^{n}$}; % annotate bundle

% % Draw fibre bundle bottom
% \draw[thick] (0,0) -- (2,0);
% \draw[thick] (2,0) -- (3,1);
% \draw[thick] (3,1) -- (1,1);
% \draw[thick] (1,1) -- (0,0);

% % Draw ray
% \draw[thick, HighlightColor] (1.25,0.75) -- (1.25,2+0.75);
% \filldraw[fill=white, draw=HighlightColor] (1.25,0.75) circle (0.04);
% \fill[HighlightColor] (1.25,2+0.75) circle (0.04);
% \node[below] at (1.25,0.75) {\footnotesize\color{white}$[\psi]$}; % annotate ray

% % Draw point on ray
% \fill (1.25,1.35) circle (0.04);
% \node[right] at (1.25,1.35) {\footnotesize$\ket{\psi}$};
% \fill (1.25,2.5) circle (0.04);
% \node[right] at (1.25,2.5) {\footnotesize$\Exp{i\phi}\!\ket{\psi}$};

% % Draw fibre bundle top
% \draw[thick] (0,0+2) -- (2,0+2);
% \draw[thick] (2,0+2) -- (3,1+2);
% \draw[thick] (3,1+2) -- (1,1+2);
% \draw[thick] (1,1+2) -- (0,0+2);

% % Draw fibre bundle sides
% \draw[thick] (0,0) -- (0,0+2);
% \draw[thick] (2,0) -- (2,0+2);
% \draw[thick] (3,1) -- (3,1+2);
% \draw[thick] (1,1) -- (1,1+2);

% Draw manifold
\fill[MainColor] (0,0) -- (2,0) -- (3,1) -- (1,1) -- cycle;
\node at (2.4,0) {\color{MainColor}$\mathscr{M}$};

% Draw curve with curve
\draw[white, thick] (0.5,0.15) to[bend left]node[midway, circle, inner sep=0, minimum size=0.1cm, fill=white, name=vectorpoint] {} (2.5,0.75);
\node at (1.25,0.25) {\footnotesize\color{white}$\ve{V}(\lambda)$}; % vector label

% Draw tangent vector
\draw[-stealth, thick, HighlightColor] ($(vectorpoint)$) -- (2,0.9);
\node at (1.75,0.6) {\footnotesize\color{HighlightColor}$\D\ve{V}$}; % vector label



\end{tikzpicture}
}}
        \visible<5->{Metric for quantum states?
        \eqn{
            Q_{\mu\nu} \overset{?}{=} \braket{\del_{\mu}\psi}{\del_{\nu}\psi}
        }}
        \visible<6->{\textbf{Close, but no cigar!}}

        \visible<7->{$\ket{\del_{\mu}\psi}$ tangent to $S^{2n-1}\neq\mathbb{CP}^{n-1}$}
    \end{columns}
\end{frame}

\begin{frame}\frametitle{Derivatives of Quantum States}
    \begin{columns}
        \column{0.6\linewidth}
        \visible<2->{Two states on same ray, careful!
        $$\ket{\psi} = \Exp{i\phi}\!\ket{\chi}~~\not\Longrightarrow~~\ket{\del\psi} \sim \ket{\del\chi}$$}
        \visible<3->{Parallel component is \textbf{not} intrinsic
        $$\ket{\partial\psi} = \Exp{i\phi}\!\ket{\partial\chi} + \alert{i\partial\phi\!\ket{\psi}}$$}
        \visible<4->{Perpendicular component \textbf{is}
        $$\ket{\mathcal{D}\psi}\equiv\p{\mathbb{1} - \ketbra{\psi}{\psi}}\!\ket{\partial\psi}$$}

        \column{0.4\linewidth}
            \visible<3->{\resizebox{\linewidth}{!}{\begin{tikzpicture}

\draw[thick, -Stealth] (0,0) -- node[midway, above, xshift=-0.1cm] {$\ket{\partial\psi}$} (4,2);

\draw[thick, -Stealth, MainColor] (0,0) -- node[midway, below] {$\Exp{i\phi}\!\ket{\partial\chi}$} (4,0);
\draw[thick, -Stealth, HighlightColor] (0,0) -- node[midway, left] {$i\partial\phi\!\ket{\psi}$} (0,2);

\fill[black] (0,0) circle (0.05cm) node [xshift=-0.25cm, yshift=-0.2cm] {$\ket{\psi}$};

\end{tikzpicture}
}}
            \visible<5->{\textbf{Covariant derivative}
            $$\cboxed{\ket{\mathcal{D}\psi} = \Exp{i\phi}\!\ket{\mathcal{D}\chi}}$$}
            \visible<6->{Tangent to $\mathbb{CP}^{n-1}$}
    \end{columns}
\end{frame}

\begin{frame}\frametitle{Metric on State Space: Fubini-Study}
\visible<2->{\textbf{Metric} in $\mathbb{CP}^{n-1}$: \textbf{Fubini-Study}
    {\large\eqn{\cboxed{Q_{\mu\nu} = \braket{\mathcal{D}_\mu\psi}{\mathcal{D}_\nu\psi} = \langle\del_\mu\psi|\p{1 - \ketbra{\psi}{\psi}}|\del_\nu\psi\rangle}}}}
\visible<3->{Manifestly gauge invariant
\eqn{
    \left\langle\mathcal{D}_\mu\psi\middle|\mathcal{D}_\nu\psi\right\rangle = \left\langle\mathcal{D}_\mu\chi\middle|\Exp{-i\phi}\Exp{i\phi}\middle|\mathcal{D}_\nu\chi\right\rangle = \left\langle\mathcal{D}_\mu\chi\middle|\mathcal{D}_\nu\chi\right\rangle
}}
\visible<4->{In litterature, separated into two parts: \textbf{\color{MainColor}quantum metric} and \textbf{\alert{Berry curvature}}
{\large$$
    Q_{\mu\nu} = {\color{MainColor}{g_{\mu\nu}}} - \frac{i}{2}\alert{\mathcal{F}_{\mu\nu}}
$$}}
\visible<5->{Since $Q_{\mu\nu}$ is \textbf{hermitian}
\eqn{
    &{\color{MainColor}g_{\mu\nu} = g_{\nu\mu}~~~(\text{symmetric})} & &{\color{HighlightColor}\mathcal{F}_{\mu\nu} = -\mathcal{F}_{\nu\mu}~~~(\text{antisymmetric})}
}}
\end{frame}

\begin{frame}\frametitle{Quantum metric}
\begin{columns}
\column{0.7\linewidth}
\visible<2->{Distance between two neigboring rays
$$\D[\psi]^2 = Q_{\mu\nu}\D{x}^{\mu}\D{x}^{\nu} = {\color{MainColor}g_{\mu\nu}}\D{x}^{\mu}\D{x}^{\nu} - \frac{i}{2}\underbrace{\cancel{{\color{HighlightColor}\mathcal{F}_{\mu\nu}}\D{x}^{\mu}\D{x}^{\nu}}}_{\text{antisymmetry}}$$}
\visible<3->{Only \textbf{{\color{MainColor}quantum metric} encodes distance} between rays!}

\column{0.3\linewidth}
\visible<3->{\resizebox{\linewidth}{!}{\begin{tikzpicture}


% Draw projective space section
\fill[MainColor] (0,0) -- (2,0) -- (3,1) -- (1,1) -- cycle;
\node[right] at (2,0) {\footnotesize\color{MainColor}$\mathbb{CP}^{n-1}$}; % annotate section

% Draw ray
\draw[thick, HighlightColor] (1.25,0.75) -- (1.25,2+0.75);
\filldraw[fill=white, draw=HighlightColor] (1.25,0.75) circle (0.04);
\fill[HighlightColor] (1.25,2+0.75) circle (0.04);

% Draw neighboring ray
\draw[thick, HighlightColor] (1.25+0.7,0.75-0.2) -- (1.25+0.7,2+0.75-0.2);
\filldraw[fill=white, draw=HighlightColor] (1.25+0.7,0.75-0.2) circle (0.04);
\fill[HighlightColor] (1.25+0.7,2+0.75-0.2) circle (0.04);

% Draw path
\draw[thick, white] (1.25,0.75) to[bend right] node [below] {\color{white}$\D[\psi]^2$} (1.25+0.7,0.75-0.2);
\foreach \i in {1,2,3,4} { % draw suggestion of rays being compared
    \draw[thick, dashed, HighlightColor] (1.25,0.75+0.5*\i) to[bend right] (1.25+0.7,0.75-0.2+0.5*\i);
}

\end{tikzpicture}
}}
\end{columns}

\visible<4->{\center{\Large Where does the \textbf{\color{HighlightColor}Berry curvature} appear?}}


\end{frame}

\begin{frame}\frametitle{Berry curvature}

\begin{columns}
	\column{0.6\linewidth}
	\visible<2->{\textbf{Intuition} from explicit form
	\eqn{
	    \mathcal{F}_{\mu\nu} &= -2\,\symup{Im}\p{\braket{\mathcal{D}_\mu\psi}{\mathcal{D}_\nu\psi}}\\
					&= i\p{\braket{\mathcal{D}_\mu\psi}{\mathcal{D}_\nu\psi} - \braket{\mathcal{D}_\nu\psi}{\mathcal{D}_\mu\psi}}\\
					&= \highlightcboxed{\bra{\psi}i\cro{\mathcal{D}_\mu,\mathcal{D}_\nu}\ket{\psi}} \visible<3->{~\longleftarrow\text{Failure to close}}
	}}
	\visible<6->{For general loop $\mathcal{C} = \del\mathcal{S}$
	\eqn{\alert{\phi} = \oint_{\mathcal{C}}\D{\alert{\phi}} = \iint_{\mathcal{S}} \alert{F_{12}}\D{x}^1\D{x}^2}}
	\column{0.4\linewidth}
	\visible<4->{Close what?} \visible<5->{Close a \textbf{loop}!}
	\visible<5->{\resizebox{\linewidth}{!}{\begin{tikzpicture}

\draw[MainColor, -stealth, thick] (0,0) -- (2,0);
\draw[MainColor, -stealth, thick] (2,0) -- node[midway, right] {\color{MainColor}$[\psi]$} (3,1);
\draw[MainColor, -stealth, thick] (3,1) -- (1,1);
\draw[MainColor, -stealth, thick] (1,1) -- (0,0);

\draw[dashed] (0,0) -- (0,2);
\draw[dashed] (2,0) -- (2,2);
\draw[dashed] (3,1) -- (3,3);
\draw[dashed] (1,1) -- (1,3);

\draw[black, -stealth, thick] (0,0) -- (2,0+0.3);
\draw[black, -stealth, thick] (2,0+0.3) -- (3,1+2*0.3);
\draw[black, -stealth, thick] (3,1+2*0.3) -- node[near start, above] {$\psi$} (1,1+3*0.3);
\draw[black, -stealth, thick] (1,1+3*0.3) -- (0,0+4*0.3);

\draw[HighlightColor, thick] (0,0) -- node[midway, left] {\color{HighlightColor}$\Exp{i\phi}$} (0,1.2);

\fill[black] (0,0) circle (0.03);
\fill[black] (0,1.2) circle (0.03);

\end{tikzpicture}
}}
	\visible<5->{\eqn{
		    \alert{\Exp{i\phi}} &\simeq 1 + iF_{12}\D{x}^{1}\D{x}^{2}
		}}

\end{columns}

\end{frame}

\begingroup\setbeamertemplate{navigation symbols}{}
\begin{frame}[noframenumbering]
	\center{
	    \LARGE\sffamily\bfseries\color{HighlightColor} Where is geometry in condensed matter?
	}
	\center{
	    \resizebox{0.2\linewidth}{!}{\begin{tikzpicture}

% Material
\fill[MainColor] (0,0) -- (2,0) -- (2,3) -- (0,3) -- cycle;

% Electric field
\draw[-stealth, white, thick] (0.1,0.2) -- (0.1,2.8) node [near end, right] {$\ve{E}$};

% Berry curvature
\draw[thick, HighlightColor] (1.5,2.5) circle (0.1);
\fill[HighlightColor] (1.5,2.5) circle (0.04);
\node[above right] at (1.5,2.5) {\color{HighlightColor}$\ve{F}$};

% Draw electron
\node at (1,1.5) {\color{white}\Large e};
\draw[-stealth, thick, white] (1.2,1.5) -- (1.8,1.5) node [near end, above] {\color{white}$\ve{v}_{\text{a}}$};

\end{tikzpicture}
}
	}
\end{frame}
\endgroup

\begin{frame}\frametitle{State Space in a Crystal}
    \begin{columns}
    \column{0.6\linewidth}
    \visible<2->{Lattice \textbf{translation symmetry}
    \eqn{
        \ket{\Psi(\ve{r}+\ve{a})} \sim \ket{\Psi(\ve{r})}
    }}
    \visible<3->{Decompose on basis of \textbf{Bloch states}
    \eqn{
        \ket{\psi_{\ve{k}}(\ve{r})} = \Exp{i\ve{k}\cdot\ve{r}}\ket{u_{\ve{k}}(\ve{r})}
    }}
    \visible<4->{\textbf{Wavevectors} $\ve{k}\in\text{1BZ}$: \textbf{parameter space}}
    \column{0.4\linewidth}
    \visible<2->{\resizebox{\linewidth}{!}{\begin{tikzpicture}

    % Lattice sites
    \foreach \i in {0,1,2} {
        \foreach \j in {0,1,2,3} {
            \fill (\i*2,\j) circle (0.1cm);
        }
    }

    % Label interesting sites
    % % position labels
    % \node[below] at (1.5,1) {$\ve{r}$};
    % \node[below] at (1.5,2) {$\ve{r}+\ve{a}$};

    % State labels
    \node[above right] at (2,1) {$\ket{\Psi(\ve{r})}$};
    \node[above right] at (2,2) {$\ket{\Psi(\ve{r}+\ve{a})}$};

    % Lattice vector
    \draw[very thick, MainColor, -stealth] (2,1) -- node [left, midway] {$\ve{a}$} (2,2);

    % similitude
    \draw[very thick, stealth-stealth, HighlightColor] (2.5,1.5) to[bend right] node [right] {$\sim$} (2.5,2.2);

\end{tikzpicture}
}}
    \end{columns}
\end{frame}

\begin{frame}\frametitle{Anomalous Velocity}

        \visible<2->{Take \textbf{semiclassical} limit (\textbf{packet})
        \eqn{
            \Psi(\ve{r}, \ve{k}) \longrightarrow \Delta{\ve{r}}, \Delta{\ve{k}}, E, \dots
        }}
        \visible<3->{\textbf{Equations of motion} for $\dot{\ve{r}}$ and $\dot{\ve{k}}$
        \eqn{
            \hbar\dot{\ve{k}} &= -e\cro{\ve{E} + \dot{\ve{r}}\times\ve{B}}\\
            \dot{\ve{r}} &= \frac{1}{\hbar}\ddp[\varepsilon]{k} - \ve{\dot{k}}\times\alert{\ve{F}}
        }}
        \visible<4->{Berry \textbf{curvature} gives \enquote{\textbf{anomalous velocity}}
        \eqn{
            {\color{HighlightColor}F_{a}} \equiv \frac{1}{2}\varepsilon_{abc}F_{bc}
        }}\vspace{-2ex}
        \visible<5->{\center{\inlinecbox{\Large\textbf{Where} can we \textbf{see this}?}}}
\end{frame}

\begin{frame}\frametitle{Chern Insulators}
    \begin{columns}
        \column{0.4\linewidth}
        \visible<2->{\center{\resizebox{0.7\linewidth}{!}{\begin{tikzpicture}

% Material
\fill[MainColor] (0,0) -- (2,0) -- (2,3) -- (0,3) -- cycle;

% Electric field
\draw[-stealth, white, thick] (0.1,0.2) -- (0.1,2.8) node [near end, right] {$\ve{E}$};

% Berry curvature
\draw[thick, HighlightColor] (1.5,2.5) circle (0.1);
\fill[HighlightColor] (1.5,2.5) circle (0.04);
\node[above right] at (1.5,2.5) {\color{HighlightColor}$\ve{F}$};

% Draw electron
\node at (1,1.5) {\color{white}\Large e};
\draw[-stealth, thick, white] (1.2,1.5) -- (1.8,1.5) node [near end, above] {\color{white}$\ve{v}_{\text{a}}$};

\end{tikzpicture}
}}
        \visible<5->{\Large\eqn{
        \sigma_{yx} = \frac{e^2}{h}C
        }}}

        \column{0.6\linewidth}
        \visible<2->{One filled band: \textbf{ignore normal term}
        \eqn{
            \ve{v}_a = \frac{e}{\hbar}\ve{E}\times\ve{F}
        }}
        \visible<3->{Calculate \textbf{Hall conductivity}
        \eqn{
            \sigma_{yx} = \frac{e^2}{h}\frac{1}{2\pi}\int_{\text{1BZ}}F_{xy}\D{k_x}\D{k_y}
        }}
        \visible<4->{Chern number is \textbf{quantized} (Gauss-Bonnet)
        \eqn{
            C = \frac{1}{2\pi}\int_{\text{1BZ}}F_{xy}\D{k_x}\D{k_y} \in \mathbb{Z}
        }}
    \end{columns}
\end{frame}

\begin{frame}\frametitle{A link between curvature and volume}
    \begin{columns}
        \column{0.6\linewidth}
        \visible<2->{\textbf{2D crystal}: Fubini-Study is simple
        \eqn{
            \symup{Q} = \pmat{{\color{MainColor}g_{xx}} & {\color{MainColor}g_{xy}} - \frac{i}{2}{\color{HighlightColor}F_{xy}}\\
            {\color{MainColor}g_{xy}} + \frac{i}{2}{\color{HighlightColor}F_{xy}} & {\color{MainColor}g_{yy}} \\}
        }}
        \visible<3->{$\Det{\symup{Q}}\geq0$ gives \textbf{local constraint}
        \eqn{
            \sqrt{\Det{\color{MainColor}\symup{g}}}\geq \frac{1}{2}\abs{\color{HighlightColor}F_{xy}}
        }}
        \visible<4->{Integration yields \textbf{global constraint}
        \eqn{
            \int_{\text{1BZ}} \sqrt{\Det{\color{MainColor}\symup{g}}}\,\D^2k &\geq \frac{1}{2}\abs{\int_{\text{1BZ}}{\color{HighlightColor}F_{xy}}\,\D^2k}\\
            {\color{MainColor}\mathcal{V}} &\geq \pi\abs{\color{HighlightColor}C}
        }}


        \column{0.4\linewidth}
        \visible<5->{\textbf{Volume} of \textbf{1BZ's image} in $\mathbb{CP}^{n-1}$
        \eqn{
            {\color{MainColor}\mathcal{V}} = \int_{\text{1BZ}} \sqrt{\Det{\color{MainColor}\symup{g}}}\,\D^2k
        }

        \center{\resizebox{0.5\linewidth}{!}{\begin{tikzpicture}

% Draw the 1BZ as a torus
\draw[thick] (0,0) ellipse (2 and 1); % outside
\draw[thick] (-1,0.1) to[bend right] (1,0.1); % bottom inner perspective
\draw[thick] (-0.8,0) to[bend left] (0.8,0); % bottom outer perspective
\node at (0,-0.55) {\Large 1BZ}; % label

% Projective space ball
\draw[thick, dashed] (0,4) circle (2cm);
\node at (0,5) {\Large$\mathbb{CP}^{n-1}$};

% Image of the 1BZ
\fill[MainColor] (0,3) circle (1cm);
\node at (0,3) {\Large\color{white}$[\psi(\ve{k})]$};

% Arrow from params to target
\draw[-stealth, very thick, HighlightColor] (-1,0.7) to[bend left] (-0.6,2.6);

\end{tikzpicture}
}}}
    \end{columns}
\end{frame}

\begin{frame}\frametitle{Geometry in Quantum Theory...}

    \vspace{-3ex}
    \begin{columns}[t]
        \column{0.33\linewidth}
        \center{\color{MainColor}\LARGE\sffamily\bfseries What}

        \visible<2->{Distinct states $[\psi]\in\mathbb{CP}^{n-1}$

        \resizebox{0.8\linewidth}{!}{\begin{tikzpicture}

\node[above] at (2,3) {\footnotesize$S^{2n-1}\subset\mathbb{C}^{n}$}; % annotate bundle

% Draw fibre bundle bottom
\draw[thick] (0,0) -- (2,0);
\draw[thick] (2,0) -- (3,1);
\draw[thick] (3,1) -- (1,1);
\draw[thick] (1,1) -- (0,0);

% Draw projective space section
\fill[MainColor] (0,0) -- (2,0) -- (3,1) -- (1,1) -- cycle;
\node[right] at (2,0) {\footnotesize\color{MainColor}$\mathbb{CP}^{n-1}$}; % annotate section

% Draw ray
\draw[thick, HighlightColor] (1.25,0.75) -- (1.25,2+0.75);
\filldraw[fill=white, draw=HighlightColor] (1.25,0.75) circle (0.04);
\fill[HighlightColor] (1.25,2+0.75) circle (0.04);
\node[below] at (1.25,0.75) {\footnotesize\color{white}$[\psi]$}; % annotate ray

% Draw point on ray
\fill (1.25,1.35) circle (0.04);
\node[right] at (1.25,1.35) {\footnotesize$\ket{\psi}$};
\fill (1.25,2.5) circle (0.04);
\node[right] at (1.25,2.5) {\footnotesize$\Exp{i\phi}\!\ket{\psi}$};

% Draw fibre bundle top
\draw[thick] (0,0+2) -- (2,0+2);
\draw[thick] (2,0+2) -- (3,1+2);
\draw[thick] (3,1+2) -- (1,1+2);
\draw[thick] (1,1+2) -- (0,0+2);

% Draw fibre bundle sides
\draw[thick] (0,0) -- (0,0+2);
\draw[thick] (2,0) -- (2,0+2);
\draw[thick] (3,1) -- (3,1+2);
\draw[thick] (1,1) -- (1,1+2);

% \draw[MainColor, -stealth, thick] (0,0) -- (2,0);
% \draw[MainColor, -stealth, thick] (2,0) -- node[midway, right] {\color{MainColor}$[\psi]$} (3,1);
% \draw[MainColor, -stealth, thick] (3,1) -- (1,1);
% \draw[MainColor, -stealth, thick] (1,1) -- (0,0);

% \draw[dashed] (0,0) -- (0,2);
% \draw[dashed] (2,0) -- (2,2);
% \draw[dashed] (3,1) -- (3,3);
% \draw[dashed] (1,1) -- (1,3);

% \draw[black, -stealth, thick] (0,0) -- (2,0+0.3);
% \draw[black, -stealth, thick] (2,0+0.3) -- (3,1+2*0.3);
% \draw[black, -stealth, thick] (3,1+2*0.3) -- node[near start, above] {$\psi$} (1,1+3*0.3);
% \draw[black, -stealth, thick] (1,1+3*0.3) -- (0,0+4*0.3);

% \draw[HighlightColor, thick] (0,0) -- node[midway, left] {\color{HighlightColor}$\Exp{i\phi}$} (0,1.2);

% \fill[black] (0,0) circle (0.03);
% \fill[black] (0,1.2) circle (0.03);

\end{tikzpicture}
}

        Remove \textbf{\color{MainColor}norm} and \textbf{\color{HighlightColor}phase}}



        \column{0.33\linewidth}
        \center{\color{MainColor}\LARGE\sffamily\bfseries How}

        \visible<3->{Fubini-Study metric
        \vspace{-2ex}
        \eqn{
            Q_{\mu\nu} &= \braket{\mathcal{D}_\mu\psi}{\mathcal{D}_\nu\psi}\\
            {\color{MainColor}g_{\mu\nu}} &= \operatorname{Re}Q_{\mu\nu}\\
            \alert{F_{\mu\nu}} &= -2\operatorname{Im}Q_{\mu\nu}
        }
        \vspace{-3ex}
        \resizebox{0.5\linewidth}{!}{\begin{tikzpicture}


% Draw projective space section
\fill[MainColor] (0,0) -- (2,0) -- (3,1) -- (1,1) -- cycle;
\node[right] at (2,0) {\footnotesize\color{MainColor}$\mathbb{CP}^{n-1}$}; % annotate section

% Draw ray
\draw[thick, HighlightColor] (1.25,0.75) -- (1.25,2+0.75);
\filldraw[fill=white, draw=HighlightColor] (1.25,0.75) circle (0.04);
\fill[HighlightColor] (1.25,2+0.75) circle (0.04);

% Draw neighboring ray
\draw[thick, HighlightColor] (1.25+0.7,0.75-0.2) -- (1.25+0.7,2+0.75-0.2);
\filldraw[fill=white, draw=HighlightColor] (1.25+0.7,0.75-0.2) circle (0.04);
\fill[HighlightColor] (1.25+0.7,2+0.75-0.2) circle (0.04);

% Draw path
\draw[thick, white] (1.25,0.75) to[bend right] node [below] {\color{white}$\D[\psi]^2$} (1.25+0.7,0.75-0.2);
\foreach \i in {1,2,3,4} { % draw suggestion of rays being compared
    \draw[thick, dashed, HighlightColor] (1.25,0.75+0.5*\i) to[bend right] (1.25+0.7,0.75-0.2+0.5*\i);
}

\end{tikzpicture}
}\resizebox{0.5\linewidth}{!}{\begin{tikzpicture}

\draw[MainColor, -stealth, thick] (0,0) -- (2,0);
\draw[MainColor, -stealth, thick] (2,0) -- node[midway, right] {\color{MainColor}$[\psi]$} (3,1);
\draw[MainColor, -stealth, thick] (3,1) -- (1,1);
\draw[MainColor, -stealth, thick] (1,1) -- (0,0);

\draw[dashed] (0,0) -- (0,2);
\draw[dashed] (2,0) -- (2,2);
\draw[dashed] (3,1) -- (3,3);
\draw[dashed] (1,1) -- (1,3);

\draw[black, -stealth, thick] (0,0) -- (2,0+0.3);
\draw[black, -stealth, thick] (2,0+0.3) -- (3,1+2*0.3);
\draw[black, -stealth, thick] (3,1+2*0.3) -- node[near start, above] {$\psi$} (1,1+3*0.3);
\draw[black, -stealth, thick] (1,1+3*0.3) -- (0,0+4*0.3);

\draw[HighlightColor, thick] (0,0) -- node[midway, left] {\color{HighlightColor}$\Exp{i\phi}$} (0,1.2);

\fill[black] (0,0) circle (0.03);
\fill[black] (0,1.2) circle (0.03);

\end{tikzpicture}
}
        \vspace{-4ex}
        \eqn{
            \D[\psi]^2 &= {\color{MainColor}g_{\mu\nu}}\D{x}^{\mu}\D{x}^{\nu}\\
            \phi &= \iint_{\mathcal{S}}{\color{HighlightColor}F_{\mu\nu}}\D{x}^{\mu}\D{x}^{\nu}
        }}

        \column{0.33\linewidth}
        \center{\color{MainColor}\LARGE\sffamily\bfseries Where}

        \visible<4->{Electron dynamics
        \resizebox{0.5\linewidth}{!}{\begin{tikzpicture}

% Material
\fill[MainColor] (0,0) -- (2,0) -- (2,3) -- (0,3) -- cycle;

% Electric field
\draw[-stealth, white, thick] (0.1,0.2) -- (0.1,2.8) node [near end, right] {$\ve{E}$};

% Berry curvature
\draw[thick, HighlightColor] (1.5,2.5) circle (0.1);
\fill[HighlightColor] (1.5,2.5) circle (0.04);
\node[above right] at (1.5,2.5) {\color{HighlightColor}$\ve{F}$};

% Draw electron
\node at (1,1.5) {\color{white}\Large e};
\draw[-stealth, thick, white] (1.2,1.5) -- (1.8,1.5) node [near end, above] {\color{white}$\ve{v}_{\text{a}}$};

\end{tikzpicture}
}
        \eqn{
            \sigma_{yx} = \frac{e^2}{h}{\color{HighlightColor}C} ~~,~~\ve{B}=\ve{0}
        }
        \eqn{
            {\color{MainColor}\mathcal{V}}\geq\pi\abs{\color{HighlightColor}C}
        }}

    \end{columns}
\end{frame}

\begingroup\setbeamertemplate{navigation symbols}{}

\begin{frame}[noframenumbering]\frametitle{Further Reading}
\visible<2->{
\textbf{Overview of topology and geometry in physics (from scratch!)}
\begin{itemize}
	\item[$\rightarrow$]\fullcite{Nakahara}
\end{itemize}
\vspace{2ex}
}

\visible<3->{
\textbf{Excellent (in house) lecture notes}
\begin{itemize}
    \item[$\rightarrow$]\fullcite{IonNotes}
\end{itemize}
}
\end{frame}



\begin{frame}[noframenumbering]
\center{
	    \huge\sffamily\bfseries\color{HighlightColor} Questions?
}
\end{frame}

\endgroup

\end{document}
